% ============================================================================
%  Laboratorio di Programmazione -- Lezione 7 (27 ottobre 2026)
%  Compilare con:  pdflatex -shell-escape Lezione07_Ereditarieta_Polimorfismo.tex
%  Richiede: mermaid.sty (nel progetto) + Mermaid CLI `mmdc` nel PATH
%  NB: ogni frame che contiene codice (lstlisting) o un diagramma (mermaid)
%      deve essere dichiarato [fragile].
% ============================================================================
\documentclass[10pt,aspectratio=169]{beamer}
\usetheme{Madrid}
\usecolortheme{default}

% `provide=*` carica la localizzazione italiana dai file .ini di babel: funziona anche
% su installazioni TeX prive di texlive-lang-italian; con una TeX Live completa e`
% equivalente al classico \usepackage[italian]{babel}.
\usepackage[italian,provide=*]{babel}
\usepackage[utf8]{inputenc}

% Localizzazione dei nomi dei blocchi Beamer
\uselanguage{italian}
\languagepath{italian}
\deftranslation[to=italian]{Definition}{Definizione}
\deftranslation[to=italian]{Theorem}{Teorema}
\deftranslation[to=italian]{Example}{Esempio}
\deftranslation[to=italian]{Proof}{Dimostrazione}
\deftranslation[to=italian]{Corollary}{Corollario}
\deftranslation[to=italian]{Lemma}{Lemma}

\usepackage{listings}
\usepackage{xcolor}
\usepackage{hyperref}
\usepackage{array}
\usepackage{booktabs}
\usepackage{mermaid}

% Mermaid: sfondo trasparente e dimensione massima adatta a una slide 16:9
\MermaidRendererOptions{-b transparent}
\MermaidAdjustBoxOpts{max width=0.95\linewidth,max height=0.62\textheight,center}

% Configurazione colori per il codice
\definecolor{codegreen}{rgb}{0,0.6,0}
\definecolor{codegray}{rgb}{0.5,0.5,0.5}
\definecolor{codepurple}{rgb}{0.58,0,0.82}
\definecolor{backcolour}{rgb}{0.95,0.95,0.92}
\definecolor{keywordblue}{rgb}{0.0,0.0,0.7}

% Stile per snippet Java
\lstdefinestyle{javastyle}{
    backgroundcolor=\color{backcolour},
    commentstyle=\color{codegreen}\itshape,
    keywordstyle=\color{keywordblue}\bfseries,
    numberstyle=\scriptsize\color{codegray},
    stringstyle=\color{codepurple},
    basicstyle=\ttfamily\footnotesize,
    breakatwhitespace=false,
    breaklines=true,
    captionpos=b,
    keepspaces=true,
    numbers=none,
    numbersep=5pt,
    showspaces=false,
    showstringspaces=false,
    showtabs=false,
    tabsize=2,
    language=Java,
    literate={à}{{\`a}}1 {è}{{\`e}}1 {é}{{\'e}}1 {ì}{{\`i}}1
             {ò}{{\`o}}1 {ù}{{\`u}}1 {È}{{\`E}}1
}
\lstset{style=javastyle}

% Stile per righe di terminale
\lstdefinestyle{shell}{
    language=bash,
    backgroundcolor=\color{black!85},
    basicstyle=\ttfamily\footnotesize\color{white},
    keywordstyle=\color{white},
    stringstyle=\color{white},
    commentstyle=\color{green!70},
    numbers=none,
    frame=none,
    showstringspaces=false
}

\title[Lezione 7 -- Ereditarietà e polimorfismo]{Lezione 7 -- Ereditarietà e polimorfismo}
\subtitle{\texttt{extends}, \texttt{super}, binding dinamico, classi astratte}
\author[Mattia Fonisto]{Mattia Fonisto}
\institute[Lab.\ di Programmazione]{Laboratorio di Programmazione -- Ingegneria Informatica, II anno\\
           Università degli Studi di Napoli Federico II -- Sede di San Giovanni a Teduccio}
\date{27 ottobre 2026}

\begin{document}

\frame{\titlepage}

% ========================================================================
\begin{frame}{Agenda della lezione (2h)}
\begin{enumerate}
    \item \textbf{Ripasso e Homework 6} \hfill (\emph{15 min})
    \begin{itemize}
        \item uno di voi alla lavagna con \texttt{Segment}, \texttt{Frazione} o \texttt{Misure}: lo commentiamo insieme
    \end{itemize}
    \item \textbf{Ereditarietà} \hfill (\emph{25 min})
    \begin{itemize}
        \item la relazione \emph{è-un}; \texttt{extends}, \texttt{super} e i costruttori
        \item che cosa vede la sottoclasse: \texttt{private} e \texttt{protected}; ridefinire un metodo, \texttt{@Override}
    \end{itemize}
    \item \textbf{Polimorfismo} \hfill (\emph{20 min})
    \begin{itemize}
        \item upcasting e binding dinamico; array polimorfici; downcasting
        \item classi e metodi astratti
    \end{itemize}
    \item \textbf{Esercitazione} -- tutti insieme \hfill (\emph{50 min})
    \begin{itemize}
        \item \emph{Commentiamo insieme}: \texttt{Forma}, \texttt{Cerchio}, \texttt{Rettangolo} -- un array di forme
        \item \emph{Insieme, passo passo}: \texttt{Dipendente} e \texttt{Manager}
    \end{itemize}
    \item \textbf{Wrap-up e Homework 7} \hfill (\emph{10 min})
\end{enumerate}
\end{frame}

% ========================================================================
\section{Ereditarietà}

\begin{frame}[fragile]{Ereditarietà: la relazione \emph{è-un}}
\begin{columns}[T]
\column{0.56\textwidth}
\begin{block}{Ereditarietà}
\small Una classe nuova, la \emph{sottoclasse}, si definisce a partire da una classe esistente, la \emph{superclasse}: ne eredita variabili di istanza e metodi, e ne aggiunge di propri o ne ridefinisce alcuni.
\end{block}

\small
È la generalizzazione-specializzazione della lezione~5: \emph{uno studente è una persona}. Tutto ciò che vale per una persona (nome, cognome, descrizione) vale per uno studente, che in più ha una matricola.

\vspace{0.1cm}
\begin{itemize}\itemsep=3pt
    \item \textbf{è-un}: un oggetto \texttt{Studente} è anche un oggetto \texttt{Persona}, e si può usare dovunque serva una \texttt{Persona}: è il polimorfismo
    \item \textbf{riuso}: \texttt{Studente} non riscrive \texttt{nome}, \texttt{cognome}, \texttt{getNome}, \texttt{getCognome}: li eredita
    \item \textbf{gerarchia}: dal generale al particolare, anche su più livelli (borsista, studente, persona)
\end{itemize}

\column{0.44\textwidth}
\MermaidAdjustBoxOpts{max width=\linewidth,max height=0.56\textheight,center}
\begin{mermaid}
classDiagram
  direction TB
  class Persona {
    -nome : String
    -cognome : String
    +Persona(String, String)
    +getNome() String
    +getCognome() String
    +descrizione() String
  }
  class Studente {
    -matricola : String
    +Studente(String, String, String)
    +getMatricola() String
    +descrizione() String
  }
  Persona <|-- Studente
\end{mermaid}
\MermaidAdjustBoxOpts{max width=0.95\linewidth,max height=0.62\textheight,center}
\footnotesize
In UML la freccia con il triangolo vuoto va dalla sottoclasse alla superclasse (lezione~5). I membri ereditati non si ripetono; \texttt{descrizione} compare in entrambe perché \texttt{Studente} la \emph{ridefinisce}.
\end{columns}
\end{frame}

% ========================================================================
\begin{frame}[fragile]{Ereditarietà in Java: \texttt{extends} -- \texttt{Persona.java} e \texttt{Studente.java}}
\begin{columns}[T]
\column{0.42\textwidth}
\begin{lstlisting}[basicstyle=\ttfamily\scriptsize]
/** Una persona: nome e cognome. */
public class Persona {

    private final String nome;
    private final String cognome;

    public Persona(String nome,
                   String cognome) {
        this.nome = nome;
        this.cognome = cognome;
    }

    public String getNome() {
        return nome;
    }

    public String getCognome() {
        return cognome;
    }

    public String descrizione() {
        return nome + " " + cognome;
    }
}
\end{lstlisting}

\column{0.58\textwidth}
\begin{lstlisting}[basicstyle=\ttfamily\scriptsize]
/** Uno studente E` UNA persona con matricola. */
public class Studente extends Persona {

    private final String matricola;

    public Studente(String nome,
                    String cognome,
                    String matricola) {
        super(nome, cognome);
        this.matricola = matricola;
    }

    public String getMatricola() {
        return matricola;
    }
\end{lstlisting}
\footnotesize
\begin{itemize}\itemsep=2pt
    \item \texttt{extends Persona}: \texttt{Studente} eredita tutto ciò che \texttt{Persona} dichiara; aggiunge \texttt{matricola} e \texttt{getMatricola}; ridefinisce \texttt{descrizione} (fra due slide)
    \item \texttt{super(nome, cognome)} chiama il costruttore di \texttt{Persona}: prossima slide. \texttt{Persona} non è stata toccata (e non è quella della lezione~6)
\end{itemize}
\end{columns}
\end{frame}

% ========================================================================
\begin{frame}[fragile]{Costruire un oggetto derivato: \texttt{super(...)} -- \texttt{ProvaPersone.java}}
\begin{columns}[T]
\column{0.5\textwidth}
\begin{lstlisting}[basicstyle=\ttfamily\scriptsize]
Persona p = new Persona("Anna", "Russo");
Studente s = new Studente("Luca", "Bianchi",
        "N86001234");

System.out.println(p.descrizione());
System.out.println(s.getNome());
System.out.println(s.getMatricola());
System.out.println(s.descrizione());

// p.getMatricola();   // ERRORE
\end{lstlisting}
\begin{lstlisting}[style=shell,basicstyle=\ttfamily\scriptsize\color{white}]
Anna Russo
Luca
N86001234
Luca Bianchi, matricola N86001234
\end{lstlisting}
\MermaidAdjustBoxOpts{max width=\linewidth,max height=0.26\textheight,center}
\begin{mermaid}
---
config:
  themeVariables:
    fontSize: 20px
  flowchart:
    nodeSpacing: 14
    rankSpacing: 30
    wrappingWidth: 600
---
flowchart LR
  s["<b>stack</b><br/>s <span style='color:#9ca3af'><i>riferimento</i></span>"] --> o["<b>heap</b> · oggetto Studente<br/>parte Persona: nome <span style='color:#9ca3af'><i>riferimento</i></span><br/>cognome <span style='color:#9ca3af'><i>riferimento</i></span><br/>parte propria: matricola <span style='color:#9ca3af'><i>riferimento</i></span>"]
  style s fill:#fef3c7,stroke:#b45309
  style o fill:#dbeafe,stroke:#1d4ed8
\end{mermaid}
\MermaidAdjustBoxOpts{max width=0.95\linewidth,max height=0.62\textheight,center}

\column{0.5\textwidth}
\footnotesize
Un oggetto \texttt{Studente} contiene \textbf{tutte} le variabili di istanza di \texttt{Persona}, più le proprie: una ``parte \texttt{Persona}'' e una parte propria. La parte \texttt{Persona} la sa inizializzare solo il costruttore di \texttt{Persona}.
\begin{itemize}\itemsep=2pt
    \item \texttt{super(nome, cognome)} lo chiama: di norma è la \textbf{prima istruzione} del costruttore; poi \texttt{Studente} inizializza la parte propria. È l'analogo di \texttt{this(...)} (lezione~6), che chiama un costruttore della \emph{stessa} classe
    \item senza \texttt{super(...)} né \texttt{this(...)}, il compilatore inserisce \texttt{super()}: se la superclasse non ha un costruttore senza argomenti, errore di compilazione
\end{itemize}
\texttt{p.getMatricola()} non compila: le persone non hanno matricola. L'ereditarietà va in una direzione sola.
\end{columns}
\end{frame}

% ========================================================================
\begin{frame}[fragile]{Che cosa vede la sottoclasse: \texttt{private} e \texttt{protected}}
\begin{columns}[T]
\column{0.5\textwidth}
{\footnotesize In \texttt{Studente} (versione 1):}
\begin{lstlisting}[basicstyle=\ttfamily\scriptsize]
    @Override
    public String descrizione() {
        // return nome + " " + cognome + ...
        // ERRORE: nome e` private in Persona
        return super.descrizione()
                + ", matricola " + matricola;
    }
\end{lstlisting}
\footnotesize
\texttt{nome} fa parte di ogni \texttt{Studente}, ma è \texttt{private} in \texttt{Persona}: visibile solo in \texttt{Persona} (lezione~5). \texttt{Studente} lo legge con i metodi \texttt{public}: \texttt{getNome()}, o \texttt{super.descrizione()}.

{\footnotesize In \texttt{Persona} e \texttt{Studente} (versione 2):}
\begin{lstlisting}[basicstyle=\ttfamily\scriptsize]
    protected final String nome;
    protected final String cognome;
\end{lstlisting}
\begin{lstlisting}[basicstyle=\ttfamily\scriptsize]
    // lecito: nome e cognome sono protected
    return nome + " " + cognome
            + ", matricola " + matricola;
\end{lstlisting}

\column{0.5\textwidth}
\footnotesize
\begin{center}
\renewcommand{\arraystretch}{1.2}
\begin{tabular}{@{}lp{5.6cm}@{}}
\toprule
\textbf{Modificatore} & \textbf{Visibile in} \\
\midrule
\texttt{private} & la classe che lo dichiara \\
\texttt{protected} & la classe, le sottoclassi, il package (lezione~9) \\
\texttt{public} & tutte le classi \\
\bottomrule
\end{tabular}
\end{center}
\begin{itemize}\itemsep=2pt
    \item \texttt{protected} apre la variabile alle sottoclassi e alle classi della stessa cartella (stesso package: lezione~9): nella versione 2 anche \texttt{ProvaPersone} vede \texttt{s.nome}; un client in un'altra cartella no
    \item l'ereditarietà \textbf{non aggira l'information hiding}: le variabili \texttt{private} di \texttt{Persona} stanno in ogni \texttt{Studente}, che però non le vede
    \item \texttt{protected} ha un costo: se la variabile non è \texttt{final}, la sottoclasse può assegnarle un valore che la superclasse non ammette. Con \texttt{private} e i \texttt{get}, le sottoclassi dipendono solo dall'interfaccia pubblica di \texttt{Persona}: se \texttt{Persona} rinomina una variabile, non cambiano
\end{itemize}
\end{columns}
\end{frame}

% ========================================================================
\begin{frame}[fragile]{Ridefinire un metodo: \emph{overriding} e \texttt{@Override}}
\begin{columns}[T]
\column{0.5\textwidth}
{\footnotesize In \texttt{Persona}:}
\begin{lstlisting}[basicstyle=\ttfamily\scriptsize]
    public String descrizione() {
        return nome + " " + cognome;
    }
\end{lstlisting}
{\footnotesize In \texttt{Studente}:}
\begin{lstlisting}[basicstyle=\ttfamily\scriptsize]
    @Override
    public String descrizione() {
        return super.descrizione()
                + ", matricola " + matricola;
    }
\end{lstlisting}
{\footnotesize In \texttt{ProvaPersone}:}
\begin{lstlisting}[basicstyle=\ttfamily\scriptsize]
System.out.println(p.descrizione());
System.out.println(s.descrizione());
\end{lstlisting}
\begin{lstlisting}[style=shell,basicstyle=\ttfamily\scriptsize\color{white}]
Anna Russo
Luca Bianchi, matricola N86001234
\end{lstlisting}

\column{0.5\textwidth}
\begin{block}{Ridefinizione (\emph{overriding})}
\small La sottoclasse riscrive un metodo della superclasse con la \textbf{stessa firma} e lo \textbf{stesso tipo di risultato}, ma un altro corpo: sugli oggetti della sottoclasse vale la nuova versione.
\end{block}
\footnotesize
\begin{itemize}\itemsep=2pt
    \item \texttt{super.descrizione()} chiama la versione di \texttt{Persona} sullo stesso oggetto: la nuova versione \textbf{estende} la vecchia. Senza \texttt{super}, \texttt{descrizione()} chiamerebbe se stessa all'infinito
    \item \texttt{@Override} è un'\emph{annotazione}: ``questo metodo ne ridefinisce uno della superclasse''. Se non è vero, per esempio per un errore di battitura (\texttt{descrizone}), il compilatore lo segnala; senza l'annotazione avremmo un metodo nuovo e nessun avviso
    \item non è l'\textbf{overloading} della lezione~3: lì stesso nome e parametri \emph{diversi} nella stessa classe; qui stessi parametri in due classi diverse
\end{itemize}
\end{columns}
\end{frame}

% ========================================================================
\begin{frame}[fragile]{Ereditarietà in Java: le regole}
\begin{columns}[T]
\column{0.62\textwidth}
\footnotesize
\textbf{Che cosa si eredita e che cosa no}
\begin{itemize}\itemsep=2pt
    \item si ereditano variabili di istanza e metodi non \texttt{private}; quelli \texttt{private} fanno parte dell'oggetto, ma la sottoclasse non li vede
    \item \textbf{non} si ereditano i costruttori: ogni classe scrive i suoi, e chiama quelli della superclasse con \texttt{super(...)}
    \item i membri \texttt{static} appartengono alla classe: si usano con il suo nome e non partecipano al binding dinamico (prossime slide)
    \item un metodo \texttt{final} non si può ridefinire; una classe \texttt{final} non si può estendere (\texttt{String} lo è). È il terzo uso di \texttt{final}, dopo le variabili locali (lezione~2) e di istanza (lezione~6): sempre ``non cambia più''
\end{itemize}
\textbf{La forma della gerarchia}
\begin{itemize}\itemsep=2pt
    \item una classe ha \textbf{una sola} superclasse diretta: niente ereditarietà multipla (le interfacce, lezione~9, ne offrono una parte)
    \item è transitiva: \texttt{Borsista extends Studente} eredita anche ciò che \texttt{Studente} ha ereditato da \texttt{Persona}
    \item una classe senza \texttt{extends} estende \texttt{Object}: ogni classe ne discende e ne eredita \texttt{equals}, \texttt{toString} e altri metodi. Per questo, come detto nella lezione~6, \texttt{equals} compila anche su una classe che non lo definisce: lezione~8
\end{itemize}

\column{0.38\textwidth}
\MermaidAdjustBoxOpts{max width=\linewidth,max height=0.62\textheight,center}
\begin{mermaid}
---
config:
  class:
    hideEmptyMembersBox: true
---
classDiagram
  direction TB
  class Object {
    +equals(Object) boolean
    +toString() String
  }
  class Persona
  class Studente
  class Borsista
  Object <|-- Persona
  Persona <|-- Studente
  Studente <|-- Borsista
\end{mermaid}
\MermaidAdjustBoxOpts{max width=0.95\linewidth,max height=0.62\textheight,center}
\footnotesize
\texttt{Borsista} sarà uno degli homework. \texttt{Object} sta in \texttt{java.lang}, come \texttt{String}.
\end{columns}
\end{frame}

% ========================================================================
\section{Polimorfismo}

\begin{frame}[fragile]{Polimorfismo: \emph{upcasting} -- \texttt{Polimorfismo.java}}
\begin{columns}[T]
\column{0.52\textwidth}
\begin{lstlisting}[basicstyle=\ttfamily\scriptsize]
Persona p = new Studente("Luca", "Bianchi",
        "N86001234");   // upcasting
System.out.println(p.getNome());
System.out.println(p.descrizione());
// p.getMatricola();   // ERRORE
\end{lstlisting}
\MermaidAdjustBoxOpts{max width=\linewidth,max height=0.3\textheight,center}
\begin{mermaid}
---
config:
  themeVariables:
    fontSize: 20px
  flowchart:
    nodeSpacing: 14
    rankSpacing: 30
    wrappingWidth: 600
---
flowchart LR
  p["<b>stack</b><br/>p <span style='color:#9ca3af'><i>riferimento</i></span><br/><i>tipo Persona</i>"] --> o["<b>heap</b> · oggetto Studente<br/>parte Persona: nome <span style='color:#9ca3af'><i>riferimento</i></span><br/>cognome <span style='color:#9ca3af'><i>riferimento</i></span><br/>parte propria: matricola <span style='color:#9ca3af'><i>riferimento</i></span>"]
  style p fill:#fef3c7,stroke:#b45309
  style o fill:#dbeafe,stroke:#1d4ed8
\end{mermaid}
\MermaidAdjustBoxOpts{max width=0.95\linewidth,max height=0.62\textheight,center}
\footnotesize
Il riferimento è di tipo \texttt{Persona}; l'oggetto a cui porta è uno \texttt{Studente}. I due tipi non coincidono, e va bene così: il riferimento copiato in \texttt{p} (lezione~6) ha un tipo più generale di quello dell'oggetto.

\column{0.48\textwidth}
\begin{block}{Upcasting}
\small Un riferimento di tipo superclasse può portare a un oggetto di una qualunque sottoclasse: uno \texttt{Studente} \emph{è una} \texttt{Persona}. La conversione è \textbf{implicita}, come \texttt{int} \textrightarrow{} \texttt{double} nella lezione~4: vince il tipo più generale.
\end{block}
\footnotesize
\begin{itemize}\itemsep=3pt
    \item quali metodi si possono chiamare su \texttt{p}? Li decide il \textbf{compilatore} dal \textbf{tipo del riferimento}: quelli di \texttt{Persona}. \texttt{p.getNome()} compila, \texttt{p.getMatricola()} no, anche se l'oggetto ce l'ha
    \item \texttt{p.descrizione()} compila: \texttt{Persona} ha \texttt{descrizione}. Ma ne esistono due versioni, di \texttt{Persona} e di \texttt{Studente}: quale viene eseguita?
    \item l'oggetto non cambia: resta uno \texttt{Studente}. Cambia solo ciò che \texttt{p} permette di vedere
\end{itemize}
\end{columns}
\end{frame}

% ========================================================================
\begin{frame}[fragile]{Polimorfismo: binding dinamico}
\begin{columns}[T]
\column{0.5\textwidth}
\begin{lstlisting}[style=shell,basicstyle=\ttfamily\scriptsize\color{white}]
Luca
Luca Bianchi, matricola N86001234
\end{lstlisting}
\footnotesize
\texttt{p.descrizione()} ha eseguito la versione di \texttt{Studente}: la classe dell'\textbf{oggetto}, non del riferimento.
\vspace{-0.1cm}
\begin{block}{Polimorfismo}
\small La stessa chiamata, \texttt{x.descrizione()}, esegue codice diverso a seconda della classe dell'oggetto a cui \texttt{x} porta in quel momento. La scelta avviene \textbf{a tempo di esecuzione} (\emph{binding dinamico}): la fa la JVM, non il compilatore.
\end{block}
\footnotesize
È il terzo principio della lezione~5, \texttt{f.area()} su cerchi e rettangoli: lo faremo nell'esercitazione. Due momenti, due controlli:
\begin{itemize}\itemsep=2pt
    \item \textbf{compilazione}: \texttt{p} è \texttt{Persona}; \texttt{Persona} ha \texttt{descrizione()}? Sì: la chiamata è lecita
    \item \textbf{esecuzione}: \texttt{p} porta a uno \texttt{Studente}: la JVM cerca \texttt{descrizione()} da \texttt{Studente}
\end{itemize}

\column{0.5\textwidth}
\MermaidAdjustBoxOpts{max width=\linewidth,max height=0.44\textheight,center}
\begin{mermaid}
---
config:
  themeVariables:
    fontSize: 20px
  flowchart:
    nodeSpacing: 14
    rankSpacing: 26
    wrappingWidth: 600
---
flowchart TB
  c["<b>p.descrizione()</b><br/>p porta a uno Studente"] --> S["classe <b>Studente</b><br/>descrizione() ✔ eseguito"]
  S -. "se non c'è, risali" .-> P["classe <b>Persona</b><br/>getNome() &nbsp; descrizione()"]
  style c fill:#fef3c7,stroke:#b45309
  style S fill:#dcfce7,stroke:#15803d
  style P fill:#dbeafe,stroke:#1d4ed8
\end{mermaid}
\MermaidAdjustBoxOpts{max width=0.95\linewidth,max height=0.62\textheight,center}
\footnotesize
La JVM parte dalla classe dell'oggetto e risale finché trova il metodo. Per \texttt{p.getNome()} in \texttt{Studente} non c'è nulla: si risale a \texttt{Persona}.

\vspace{0.05cm}
Vale per i metodi che si possono ridefinire. Per i metodi \texttt{static}, \texttt{private} e \texttt{final}, per i costruttori e per la scelta fra metodi in overloading (lezione~3) decide il compilatore: \emph{binding statico}.
\end{columns}
\end{frame}

% ========================================================================
\begin{frame}[fragile]{Polimorfismo: un array di \texttt{Persona} -- \texttt{Polimorfismo.java}}
\begin{columns}[T]
\column{0.54\textwidth}
\begin{lstlisting}[basicstyle=\ttfamily\scriptsize]
/** LEGGE la persona ricevuta e la stampa. */
public static void stampa(Persona x) {
    System.out.println(x.descrizione());
}
\end{lstlisting}
{\footnotesize Nel \texttt{main}:}
\begin{lstlisting}[basicstyle=\ttfamily\scriptsize]
Persona[] elenco = {
    new Persona("Anna", "Russo"),
    new Studente("Luca", "Bianchi", "N86001234"),
    new Persona("Marco", "Esposito"),
    new Studente("Sara", "Greco", "N86005678")
};
for (Persona x : elenco) {
    stampa(x);
}
\end{lstlisting}
\begin{lstlisting}[style=shell,basicstyle=\ttfamily\scriptsize\color{white}]
Anna Russo
Luca Bianchi, matricola N86001234
Marco Esposito
Sara Greco, matricola N86005678
\end{lstlisting}

\column{0.46\textwidth}
\footnotesize
\begin{itemize}\itemsep=2pt
    \item un \texttt{Persona[]} può contenere riferimenti a persone e a studenti: ogni elemento è un upcasting. È un \textbf{array polimorfico}
    \item il ciclo e \texttt{stampa} sono scritti una volta sola, in termini di \texttt{Persona}; a ogni giro \texttt{x.descrizione()} esegue la versione giusta per l'oggetto di quel giro
    \item anche passare \texttt{x} a \texttt{stampa} è un upcasting: accetta una \texttt{Persona}, quindi anche uno \texttt{Studente}
\end{itemize}
\vspace{0.05cm}
\textbf{Che cosa si guadagna.} Se domani aggiungiamo \texttt{Docente extends Persona}, né \texttt{stampa} né il ciclo cambiano: funzionano anche per i docenti. Il codice che usa la gerarchia è scritto per la superclasse e non conosce l'elenco delle sottoclassi.
\end{columns}
\end{frame}

% ========================================================================
\begin{frame}[fragile]{Polimorfismo: \emph{downcasting} -- \texttt{Downcast.java}}
\begin{columns}[T]
\column{0.52\textwidth}
\begin{lstlisting}[basicstyle=\ttfamily\scriptsize]
Persona p = new Studente("Luca", "Bianchi",
        "N86001234");
Persona q = new Persona("Anna", "Russo");

// Studente s = p;   // ERRORE: serve il cast
Studente s = (Studente) p;   // ok: Studente
System.out.println(s.getMatricola());

Studente t = (Studente) q;   // compila, ma...
System.out.println(t.getMatricola());
\end{lstlisting}
\begin{lstlisting}[style=shell,basicstyle=\ttfamily\scriptsize\color{white}]
N86001234
Exception in thread "main"
  java.lang.ClassCastException: class Persona
  cannot be cast to class Studente ...
\end{lstlisting}
\vspace{-0.1cm}
\footnotesize
Il programma termina su \texttt{(Studente) q}: la JVM ha controllato che l'oggetto di \texttt{q} non è uno \texttt{Studente}. Il messaggio è un'\emph{eccezione}, lezione~10.

\column{0.46\textwidth}
\begin{block}{Downcasting}
\small Assegnare a un riferimento di tipo sottoclasse un riferimento di tipo superclasse. Serve il \textbf{cast} \texttt{(Studente) p}: compila perché \texttt{Studente} è sottoclasse di \texttt{Persona}, ma riesce solo se l'oggetto è uno \texttt{Studente}.
\end{block}
\begin{center}
\scriptsize
\renewcommand{\arraystretch}{1.15}
\begin{tabular}{@{}lll@{}}
\toprule
 & \textbf{Primitivi (L4)} & \textbf{Classi} \\
\midrule
implicito & \texttt{double d = i;} & \texttt{Persona p = s;} \\
esplicito & \texttt{int i = (int) d;} & \texttt{Studente s = (Studente) p;} \\
rischio & perde i decimali & errore a tempo di esecuzione \\
\bottomrule
\end{tabular}
\end{center}
\footnotesize
\begin{itemize}\itemsep=2pt
    \item per sapere \emph{prima} se l'oggetto è uno \texttt{Studente}, Java ha un operatore apposta (lezione~8): fino ad allora, il cast si fa solo se si è certi della classe dell'oggetto
\end{itemize}
\end{columns}
\end{frame}

% ========================================================================
\begin{frame}[fragile]{Classi e metodi astratti}
\begin{columns}[T]
\column{0.5\textwidth}
{\footnotesize In \texttt{Forma}, la superclasse dell'esercitazione:}
\begin{lstlisting}[basicstyle=\ttfamily\scriptsize]
public abstract class Forma {
    private final String colore;

    public Forma(String colore) {
        this.colore = colore;
    }

    public String getColore() {
        return colore;
    }

    public abstract double area();

    public boolean piuGrandeDi(Forma altra) {
        return area() > altra.area();
    }
\end{lstlisting}
{\footnotesize In \texttt{ProvaForme}: riferimenti e array di tipo \texttt{Forma} sono leciti}
\begin{lstlisting}[basicstyle=\ttfamily\scriptsize]
// Forma f = new Forma("nero");   // ERRORE
Forma[] forme = {
    new Cerchio("rosso", 1),
\end{lstlisting}

\column{0.5\textwidth}
\begin{block}{Classe astratta, metodo astratto}
\small Un metodo \texttt{abstract} ha la firma ma non il corpo: dice che cosa ogni sottoclasse deve saper fare. Una classe con un metodo astratto è \texttt{abstract}: si può estendere e usare come tipo di riferimenti, ma \textbf{non creare oggetti}.
\end{block}
\footnotesize
\begin{itemize}\itemsep=2pt
    \item una sottoclasse \textbf{deve} ridefinire tutti i metodi astratti che eredita, altrimenti è a sua volta \texttt{abstract}
    \item \texttt{piuGrandeDi} ha il corpo e chiama \texttt{area()} su \texttt{this} e su \texttt{altra}: a tempo di esecuzione gli oggetti sono cerchi e rettangoli, e la JVM esegue la loro \texttt{area}. Un metodo della superclasse può usare un metodo che ha corpo solo nelle sottoclassi
    \item \texttt{Forma} ha anche la parte comune: \texttt{colore}, il costruttore (chiamato dalle sottoclassi con \texttt{super(colore)}), \texttt{getColore}. In UML i nomi astratti sono in \emph{corsivo}
\end{itemize}
\end{columns}
\end{frame}

% ========================================================================
\section{Esercitazione}

\begin{frame}[fragile]{Commentiamo insieme (20 min) -- il problema}
\begin{columns}[T]
\column{0.55\textwidth}
\footnotesize
\textbf{Il problema.} Un programma di disegno tiene in un array le forme sul foglio: cerchi e rettangoli, ciascuno con un colore. Deve stampare la descrizione di ogni forma, calcolare l'area totale e trovare la forma più grande.

\vspace{0.1cm}
\textbf{Le regole}
\begin{itemize}\itemsep=2pt
    \item ogni forma ha un colore, un'area, un perimetro e una descrizione
    \item non esiste una forma ``generica'': sul foglio ci sono solo cerchi e rettangoli. Domani, forse, anche triangoli
    \item un cerchio ha un raggio, un rettangolo una base e un'altezza: un valore non positivo diventa 1. Dopo la creazione una forma non cambia più
    \item un rettangolo sa dire se è un quadrato
    \item di due forme si può dire quale ha l'area maggiore, qualunque forme siano
\end{itemize}

\vspace{0.1cm}
La soluzione è già scritta: la classe astratta \texttt{Forma}, le sottoclassi \texttt{Cerchio} e \texttt{Rettangolo} e il programma \texttt{ProvaForme}. Li leggiamo insieme, riga per riga, cercando \texttt{extends}, \texttt{super}, \texttt{@Override} e \texttt{abstract}.

\column{0.45\textwidth}
\MermaidAdjustBoxOpts{max width=\linewidth,max height=0.76\textheight,center}
\begin{mermaid}
classDiagram
  direction TB
  class Forma {
    <<abstract>>
    -colore : String
    +Forma(String)
    +getColore() String
    +area()* double
    +perimetro()* double
    +piuGrandeDi(Forma) boolean
    +descrizione() String
  }
  class Cerchio {
    -raggio : double
    +Cerchio(String, double)
    +getRaggio() double
    +area() double
    +perimetro() double
    +descrizione() String
  }
  class Rettangolo {
    -base : double
    -altezza : double
    +Rettangolo(String, double, double)
    +getBase() double
    +getAltezza() double
    +isQuadrato() boolean
    +area() double
    +perimetro() double
    +descrizione() String
  }
  Forma <|-- Cerchio
  Forma <|-- Rettangolo
\end{mermaid}
\MermaidAdjustBoxOpts{max width=0.95\linewidth,max height=0.62\textheight,center}
\end{columns}
\end{frame}

% ========================================================================
\begin{frame}[fragile]{Commentiamo insieme -- \texttt{Forma.java}}
\begin{columns}[T]
\column{0.5\textwidth}
\begin{lstlisting}[basicstyle=\ttfamily\scriptsize]
/** Una forma piana, con un colore. Astratta:
 *  esistono solo cerchi, rettangoli... */
public abstract class Forma {

    private final String colore;

    public Forma(String colore) {
        this.colore = colore;
    }

    public String getColore() {
        return colore;
    }
\end{lstlisting}
\footnotesize
\begin{itemize}\itemsep=2pt
    \item \texttt{colore} è comune a tutte le forme: sta qui, una volta sola, \texttt{private final}
    \item il costruttore non si chiama mai con \texttt{new Forma(...)}: lo chiamano \texttt{Cerchio} e \texttt{Rettangolo} con \texttt{super(colore)}
\end{itemize}

\column{0.5\textwidth}
\begin{lstlisting}[basicstyle=\ttfamily\scriptsize]
    /** Ogni forma la calcola a modo suo. */
    public abstract double area();

    public abstract double perimetro();

    /** true se ha area maggiore di altra. */
    public boolean piuGrandeDi(Forma altra) {
        return area() > altra.area();
    }

    public String descrizione() {
        return "forma " + colore;
    }
}
\end{lstlisting}
\footnotesize
\begin{itemize}\itemsep=2pt
    \item \texttt{area} e \texttt{perimetro} terminano con \texttt{;}: niente corpo. Ogni sottoclasse deve scriverli
    \item \texttt{piuGrandeDi} e \texttt{descrizione} hanno il corpo e si ereditano: \texttt{descrizione} è minima e le sottoclassi la ridefiniranno
\end{itemize}
\end{columns}
\end{frame}

% ========================================================================
\begin{frame}[fragile]{Commentiamo insieme -- \texttt{Cerchio.java} e \texttt{Rettangolo.java}}
\begin{columns}[T]
\column{0.5\textwidth}
\begin{lstlisting}[basicstyle=\ttfamily\scriptsize]
public class Cerchio extends Forma {

    private final double raggio;

    public Cerchio(String colore,
                   double raggio) {
        super(colore);
        if (raggio > 0) {
            this.raggio = raggio;
        } else {
            this.raggio = 1;
        }
    }

    @Override
    public double area() {
        return Math.PI * raggio * raggio;
    }

    @Override
    public String descrizione() {
        return "cerchio " + getColore()
                + ", raggio " + raggio;
    }
\end{lstlisting}

\column{0.5\textwidth}
\begin{lstlisting}[basicstyle=\ttfamily\scriptsize]
public class Rettangolo extends Forma {

    private final double base;
    private final double altezza;

    public boolean isQuadrato() {
        return base == altezza;
    }

    @Override
    public double area() {
        return base * altezza;
    }
\end{lstlisting}
\footnotesize
\begin{itemize}\itemsep=2pt
    \item \texttt{super(colore)} per primo; poi il controllo sul raggio (non positivo: 1), come in \texttt{Circle} della lezione~6. Il costruttore di \texttt{Rettangolo}, non mostrato, controlla allo stesso modo \texttt{base} e \texttt{altezza}
    \item \texttt{area} e \texttt{perimetro} danno corpo ai metodi astratti: senza uno dei due, \texttt{Cerchio} non compila (dovrebbe essere \texttt{abstract}). \texttt{descrizione} usa \texttt{getColore()}: \texttt{colore} è \texttt{private}. \texttt{isQuadrato} esiste solo in \texttt{Rettangolo}
\end{itemize}
\end{columns}
\end{frame}

% ========================================================================
\begin{frame}[fragile]{Commentiamo insieme -- \texttt{ProvaForme.java} (1/2)}
\begin{columns}[T]
\column{0.52\textwidth}
\begin{lstlisting}[basicstyle=\ttfamily\scriptsize]
/** LEGGE l'array e somma le aree. */
public static double areaTotale(Forma[] forme) {
    double totale = 0;
    for (Forma f : forme) {
        totale += f.area();
    }
    return totale;
}

/** LEGGE l'array e restituisce la forma
 *  piu` grande (precondizione: almeno una). */
public static Forma piuGrande(Forma[] forme) {
    Forma max = forme[0];
    for (Forma f : forme) {
        if (f.piuGrandeDi(max)) {
            max = f;
        }
    }
    return max;
}
\end{lstlisting}

\column{0.48\textwidth}
\small
\textbf{Da notare, leggendo:}
\begin{itemize}\itemsep=4pt
    \item \texttt{f} è un riferimento di tipo \texttt{Forma}; a ogni giro l'oggetto è un \texttt{Cerchio} o un \texttt{Rettangolo}. \texttt{f.area()} compila perché \texttt{Forma} dichiara \texttt{area}; esegue la versione della classe dell'oggetto
    \item \texttt{areaTotale} non nomina mai \texttt{Cerchio} né \texttt{Rettangolo}: funzionerà anche per il \texttt{Triangolo} dell'homework, senza cambiare una riga
    \item \texttt{piuGrandeDi} è scritto in \texttt{Forma}, ma quando viene eseguito \texttt{this} e \texttt{altra} portano a oggetti concreti: le due chiamate ad \texttt{area()} sono entrambe polimorfiche
    \item \texttt{piuGrande} restituisce un \texttt{Forma}: il client che la riceve non sa se è un cerchio o un rettangolo, e per stamparla non ha bisogno di saperlo
\end{itemize}
\end{columns}
\end{frame}

% ========================================================================
\begin{frame}[fragile]{Commentiamo insieme -- \texttt{ProvaForme.java} (2/2)}
\begin{columns}[T]
\column{0.56\textwidth}
\begin{lstlisting}[basicstyle=\ttfamily\scriptsize]
public static void main(String[] args) {
    Forma[] forme = {
        new Cerchio("rosso", 1),
        new Rettangolo("blu", 3, 4),
        new Cerchio("verde", 2),
        new Rettangolo("giallo", 2, 2)
    };
    for (Forma f : forme) {
        System.out.printf("%s: area %.2f%n",
                f.descrizione(), f.area());
    }
    System.out.printf("area totale: %.2f%n",
            areaTotale(forme));
    Forma g = piuGrande(forme);
    System.out.println(g.descrizione());

    // Forma f = new Forma("nero");   // ERRORE
    // Rettangolo r = forme[1];   // ERRORE: cast
    Rettangolo r = (Rettangolo) forme[1];
    System.out.println(r.isQuadrato());
}
\end{lstlisting}

\column{0.44\textwidth}
\begin{lstlisting}[style=shell,basicstyle=\ttfamily\scriptsize\color{white}]
cerchio rosso, raggio 1.0: area 3.14
rettangolo blu 3.0 x 4.0: area 12.00
cerchio verde, raggio 2.0: area 12.57
rettangolo giallo 2.0 x 2.0: area 4.00
area totale: 31.71
cerchio verde, raggio 2.0
false
\end{lstlisting}
\footnotesize
\textbf{Rispondiamo insieme:}
\begin{enumerate}\itemsep=1pt
    \item perché \texttt{f.descrizione()} non stampa ``forma rosso''?
    \item quante versioni di \texttt{area} esistono? Quante di \texttt{piuGrandeDi}?
    \item senza \texttt{perimetro} in \texttt{Cerchio}, quale file non compila?
    \item perché \texttt{r.isQuadrato()} vuole il cast, e \texttt{f.area()} no?
\end{enumerate}
\end{columns}
\end{frame}

% ========================================================================
\begin{frame}[fragile]{Insieme, passo passo (30 min) -- il problema}
\begin{columns}[T]
\column{0.64\textwidth}
\small
\textbf{Il problema.} Un'azienda tiene il personale in un array e deve calcolare il costo mensile degli stipendi. Ci sono dipendenti e manager: un manager è un dipendente con in più un bonus mensile.

\vspace{0.05cm}
\textbf{Le regole}
\begin{itemize}\itemsep=2pt
    \item ogni dipendente ha un nome e uno stipendio base mensile, positivo (precondizione); lo stipendio di un dipendente è la base
    \item un aumento del $p\%$ fa crescere la base del $p\%$; il nome non cambia mai
    \item un manager ha anche un bonus mensile fisso: il suo stipendio è base più bonus. L'aumento riguarda solo la base
    \item di ogni dipendente si può avere una descrizione: nome e stipendio; per un manager anche il bonus
    \item il costo mensile dell'azienda è la somma degli stipendi di tutto il personale
\end{itemize}

\vspace{0.05cm}
\texttt{ProvaDipendenti} è già scritto: mancano \texttt{Dipendente} e \texttt{Manager}. Scriviamo prima la superclasse, poi la sottoclasse.

\column{0.36\textwidth}
\MermaidAdjustBoxOpts{max width=\linewidth,max height=0.46\textheight,center}
\begin{mermaid}
---
config:
  themeVariables:
    fontSize: 20px
  flowchart:
    nodeSpacing: 14
    rankSpacing: 34
    wrappingWidth: 600
---
flowchart TB
  A["<b>ProvaDipendenti</b><br/>già scritta"] -- "usa" --> B["<b>Dipendente</b><br/>da scrivere"]
  A -- "usa" --> C["<b>Manager</b><br/>da scrivere"]
  C -- "estende" --> B
  style A fill:#fef3c7,stroke:#b45309
  style B fill:#fee2e2,stroke:#b91c1c,stroke-dasharray: 5 5
  style C fill:#fee2e2,stroke:#b91c1c,stroke-dasharray: 5 5
\end{mermaid}
\MermaidAdjustBoxOpts{max width=0.95\linewidth,max height=0.62\textheight,center}
\footnotesize
\texttt{Manager} avrà due ruoli: sarà usata da \texttt{ProvaDipendenti} e sarà sottoclasse di \texttt{Dipendente}.
\end{columns}
\end{frame}

% ========================================================================
\begin{frame}[fragile]{Insieme, passo passo -- \texttt{ProvaDipendenti.java}}
\begin{columns}[T]
\column{0.56\textwidth}
\begin{lstlisting}[basicstyle=\ttfamily\scriptsize]
public static void main(String[] args) {
    Dipendente a = new Dipendente("Anna", 1800);
    Manager m = new Manager("Luca", 2500, 700);
    Dipendente[] personale = {
        a, m, new Dipendente("Sara", 1600)
    };
    System.out.println(a.descrizione());
    System.out.println(m.descrizione());
    System.out.println(m.getBonus());
    System.out.println(costoMensile(personale));

    for (Dipendente d : personale) {
        d.aumenta(10);
    }
    System.out.println(m.descrizione());
    System.out.println(costoMensile(personale));

    Dipendente d = m;   // upcasting
    System.out.println(d.stipendio());
    // d.getBonus();   // ERRORE: d e` un Dipendente
}
\end{lstlisting}

\column{0.44\textwidth}
{\footnotesize Deve stampare:}
\begin{lstlisting}[style=shell,basicstyle=\ttfamily\scriptsize\color{white}]
Anna, stipendio 1800.0
Luca, stipendio 3200.0, bonus 700.0
700.0
6600.0
Luca, stipendio 3450.0, bonus 700.0
7190.0
3450.0
\end{lstlisting}
\footnotesize
Come nelle lezioni 5 e 6, il client viene prima delle classi: leggendolo ricaviamo l'\textbf{interfaccia} di \texttt{Dipendente} e di \texttt{Manager}. \texttt{costoMensile} è un metodo \texttt{static} di \texttt{ProvaDipendenti}, già scritto come \texttt{areaTotale}: somma \texttt{d.stipendio()}.

\vspace{0.05cm}
Due indizi: \texttt{aumenta(10)} è chiamato su tutti, manager compreso, e il bonus non cambia; \texttt{d.stipendio()} stampa 3450, non 2750.
\end{columns}
\end{frame}

% ========================================================================
\begin{frame}[fragile]{Insieme, passo passo -- dal programma di prova alle classi}
\begin{columns}[T]
\column{0.6\textwidth}
\footnotesize
\textbf{Prima di scrivere codice, decidiamo insieme:}
\begin{enumerate}\itemsep=2pt
    \item chi è la superclasse? \textrightarrow{} \texttt{Dipendente}: un manager \emph{è un} dipendente, non il contrario
    \item quali \textbf{variabili di istanza} in \texttt{Dipendente}? \textrightarrow{} \texttt{nome}, \texttt{final}, e \texttt{stipendioBase}, che cambia con \texttt{aumenta}: non \texttt{final}. \texttt{private} entrambe
    \item e in \texttt{Manager}? \textrightarrow{} solo \texttt{bonus}, \texttt{private final}: il resto è ereditato
    \item il \textbf{costruttore} di \texttt{Manager}? \textrightarrow{} tre parametri: i primi due a \texttt{super(...)}, il terzo a \texttt{bonus}
    \item che cosa ridefinisce \texttt{Manager}? \textrightarrow{} \texttt{stipendio}: \texttt{super.stipendio() + bonus}. E \texttt{descrizione}, che aggiunge il bonus a quella di \texttt{Dipendente}
    \item \texttt{aumenta} va ridefinito? \textrightarrow{} no: agisce sulla base, che è ereditata
    \item serve \texttt{protected} perché \texttt{Manager} legga la base? \textrightarrow{} no: basta \texttt{super.stipendio()}, o \texttt{getStipendioBase()}
\end{enumerate}

\column{0.4\textwidth}
\MermaidAdjustBoxOpts{max width=\linewidth,max height=0.74\textheight,center}
\begin{mermaid}
classDiagram
  direction TB
  class Dipendente {
    -nome : String
    -stipendioBase : double
    +Dipendente(String, double)
    +getNome() String
    +getStipendioBase() double
    +stipendio() double
    +aumenta(double)
    +descrizione() String
  }
  class Manager {
    -bonus : double
    +Manager(String, double, double)
    +getBonus() double
    +stipendio() double
    +descrizione() String
  }
  Dipendente <|-- Manager
\end{mermaid}
\MermaidAdjustBoxOpts{max width=0.95\linewidth,max height=0.62\textheight,center}
\end{columns}
\end{frame}

% ========================================================================
\begin{frame}[fragile]{Una possibile soluzione (1/2) -- \texttt{Dipendente.java}}
\begin{columns}[T]
\column{0.5\textwidth}
\begin{lstlisting}[basicstyle=\ttfamily\scriptsize]
/** Un dipendente: nome e stipendio base. */
public class Dipendente {

    private final String nome;
    private double stipendioBase;   // mensile

    /** Precondizione: stipendioBase > 0. */
    public Dipendente(String nome,
                      double stipendioBase) {
        this.nome = nome;
        this.stipendioBase = stipendioBase;
    }

    public String getNome() {
        return nome;
    }

    public double getStipendioBase() {
        return stipendioBase;
    }
\end{lstlisting}

\column{0.5\textwidth}
\begin{lstlisting}[basicstyle=\ttfamily\scriptsize]
    /** Lo stipendio mensile: la base. */
    public double stipendio() {
        return stipendioBase;
    }

    /** MODIFICA la base: +percentuale %. */
    public void aumenta(double percentuale) {
        stipendioBase += stipendioBase
                * percentuale / 100;
    }

    public String descrizione() {
        return nome + ", stipendio "
                + stipendio();
    }
}
\end{lstlisting}
\footnotesize
\begin{itemize}\itemsep=2pt
    \item una classe come quelle delle lezioni 5 e 6: non sa che avrà una sottoclasse. Non è astratta: i dipendenti esistono
    \item \texttt{descrizione} chiama \texttt{stipendio()} invece di leggere \texttt{stipendioBase}: è ciò che farà funzionare la descrizione di un manager
\end{itemize}
\end{columns}
\end{frame}

% ========================================================================
\begin{frame}[fragile]{Una possibile soluzione (2/2) -- \texttt{Manager.java}}
\begin{columns}[T]
\column{0.5\textwidth}
\begin{lstlisting}[basicstyle=\ttfamily\scriptsize]
public class Manager extends Dipendente {
    private final double bonus;   // mensile
    public Manager(String nome,
                   double stipendioBase,
                   double bonus) {
        super(nome, stipendioBase);
        this.bonus = bonus;
    }

    public double getBonus() {
        return bonus;
    }

    @Override
    public double stipendio() {
        return super.stipendio() + bonus;
    }

    @Override
    public String descrizione() {
        return super.descrizione()
                + ", bonus " + bonus;
    }
}
\end{lstlisting}

\column{0.5\textwidth}
\small
\begin{itemize}\itemsep=3pt
    \item \texttt{super(nome, stipendioBase)} per primo: la parte \texttt{Dipendente} la inizializza \texttt{Dipendente}
    \item \texttt{stipendio} ridefinito \textbf{estende} quello di \texttt{Dipendente}: \texttt{super.stipendio()} è la base, letta senza \texttt{protected}
    \item \texttt{m.descrizione()} esegue la versione di \texttt{Manager}, che chiama \texttt{super.descrizione()}, che chiama \texttt{stipendio()}: su un manager è quello di \texttt{Manager}, 3200 e non 2500. Il binding dinamico vale anche dentro la superclasse
    \item \texttt{d.aumenta(10)}: il metodo ereditato lavora sulla parte \texttt{Dipendente}; \texttt{bonus} non è toccato
    \item \texttt{d = m} è un upcasting: \texttt{d.stipendio()} stampa 3450. Il tipo di \texttt{d} decide che cosa compila (\texttt{d.getBonus()} no), l'oggetto che cosa è eseguito
\end{itemize}
\end{columns}
\end{frame}

% ========================================================================
\section{Wrap-up}

\begin{frame}{Riepilogo: cosa abbiamo trattato}
\footnotesize
\begin{center}
\renewcommand{\arraystretch}{1.25}
\begin{tabular}{|l|l|}
\hline
\textbf{Concetto (unità)} & \textbf{Dove compare oggi} \\
\hline
Ereditarietà, relazione \emph{è-un} (U13) & \texttt{Studente}, \texttt{Cerchio}, \texttt{Rettangolo}, \texttt{Manager} e le loro superclassi \\
\hline
\texttt{extends}, costruttori e \texttt{super(...)} (U13) & tutte le sottoclassi di oggi \\
\hline
\texttt{private} e \texttt{protected} (U13) & \texttt{Persona} versioni 1 e 2 \\
\hline
Ridefinizione, \texttt{super.metodo()}, \texttt{@Override} (U13) & \texttt{descrizione} in \texttt{Studente}; \texttt{stipendio} in \texttt{Manager} \\
\hline
Che cosa non si eredita; \texttt{final}; \texttt{Object} (U13) & slide ``le regole'' \\
\hline
Upcasting e downcasting (U14) & \texttt{Polimorfismo}; \texttt{Downcast}; \texttt{(Rettangolo) forme[1]}; \texttt{d = m} \\
\hline
Binding dinamico (U14) & \texttt{p.descrizione()}; \texttt{f.area()}; \texttt{d.stipendio()} \\
\hline
Array polimorfico (U14) & \texttt{Persona[]}; \texttt{Forma[]}; \texttt{Dipendente[]} \\
\hline
Classi e metodi astratti (U14) & \texttt{Forma}, \texttt{area()} e \texttt{perimetro()} \\
\hline
\end{tabular}
\end{center}

\vspace{0.1cm}
\textbf{Prossima lezione (L8)}: la superclasse \texttt{Object} -- \texttt{toString}, \texttt{equals(Object)} e \texttt{hashCode}, \texttt{instanceof}, \texttt{clone}; le relazioni tra classi: aggregazione e composizione, ereditarietà o contenimento?
\end{frame}

% ========================================================================
\begin{frame}{Homework 7 (da svolgere entro la prossima lezione)}
\footnotesize\setlength{\parskip}{0pt}
Tre sottoclassi, ciascuna con il suo programma di prova (\texttt{ProvaTriangolo}, \texttt{ProvaContoRisparmio}, \texttt{ProvaBorsisti}); le superclassi, già scritte, vanno copiate in \texttt{homework/} senza modifiche. Niente \texttt{instanceof} né \texttt{toString}.
\begin{enumerate}\itemsep=1pt
    \item \textbf{\texttt{Triangolo.java}}: una nuova forma per la gerarchia di oggi (copiate \texttt{Forma}, \texttt{Cerchio} e \texttt{Rettangolo}). Costruttore \texttt{Triangolo(String colore, double a, double b, double c)}, con i tre lati: se non formano un triangolo (uno non positivo, o non minore della somma degli altri due) diventano 3, 4, 5; il controllo è un metodo \texttt{private static}. \texttt{getA()}, \texttt{getB()}, \texttt{getC()}, \texttt{isEquilatero()}, \texttt{perimetro()}, \texttt{area()} con la formula di Erone $\sqrt{p(p-a)(p-b)(p-c)}$ ($p$ semiperimetro), \texttt{descrizione()}. Il programma di prova mette in un \texttt{Forma[]} tre triangoli (uno non valido), un cerchio e un rettangolo; stampa descrizione, area e perimetro di ciascuno e l'area totale (con l'\texttt{areaTotale} di oggi); infine, con un cast, dice se il secondo triangolo è equilatero.
    \item \textbf{\texttt{ContoRisparmio.java}}: estende il \texttt{ContoBancario} della lezione~6 (copiatelo). Ha un tasso annuo in percentuale e la costante di classe \texttt{SALDO\_MINIMO}, 100. Costruttore \texttt{ContoRisparmio(String intestatario, double primo, double tasso)}; \texttt{getTasso()}; \texttt{preleva} ridefinito: \texttt{false} anche se il prelievo porterebbe il saldo sotto \texttt{SALDO\_MINIMO}, altrimenti come in \texttt{ContoBancario}; \texttt{aggiungiInteressi()} versa gli interessi di un anno (saldo per tasso diviso 100). Il programma di prova mette in un \texttt{ContoBancario[]} un conto normale e uno di risparmio, entrambi con 500; chiede a ciascuno un prelievo di 450 e poi uno di 300, stampando risultati e saldi; aggiunge gli interessi al conto di risparmio; prova un trasferimento per verso.
    \item \textbf{\texttt{Borsista.java}}: estende lo \texttt{Studente} di oggi (copiate \texttt{Persona} e \texttt{Studente}). Importo annuo della borsa, \texttt{double} positivo; costruttore con quattro parametri; \texttt{getImportoBorsa()}; \texttt{rataMensile()}, l'importo diviso 12; \texttt{descrizione()} ridefinita. Il programma di prova crea un borsista e chiama su di lui un metodo di \texttt{Persona}, uno di \texttt{Studente} e uno di \texttt{Borsista}; mette in un \texttt{Persona[]} una persona, uno studente e il borsista e stampa le descrizioni con \texttt{stampa(Persona x)}; infine assegna il borsista a un riferimento \texttt{Studente} e ne stampa la descrizione.
\end{enumerate}
\end{frame}

% ========================================================================
\begin{frame}
\begin{center}
{\Huge Domande?}\\[2em]
{\normalsize \href{mailto:mattia.fonisto@unina.it}{\texttt{mattia.fonisto@unina.it}}}
\end{center}
\end{frame}

\end{document}
